\documentclass{article}

\usepackage{amsmath}
\usepackage{amssymb}
\usepackage{enumitem}
\usepackage[polish]{babel}
\usepackage[utf8]{inputenc}
\usepackage[T1]{fontenc}

\title{Jak sprawdzać liniową niezależność?}
\date{\today}

\begin{document}
\maketitle

\section{Definicja liniowej niezależności}

Niech $V$ będzie przestrzenią liniową nad ciałem $\mathbb K$, gdzie $\mathbb K=\mathbb R$ albo $\mathbb K=\mathbb C$. Elementy
\[
v_1,v_2,\ldots,v_n\in V
\]
są liniowo niezależne, jeżeli
\[
\alpha_1v_1+\alpha_2v_2+\cdots+\alpha_nv_n=0
\quad\Longrightarrow\quad
\alpha_1=\alpha_2=\cdots=\alpha_n=0.
\]

Jeżeli istnieją skalary $\alpha_i$, nie wszystkie równocześnie zerowe, dla których
\[
\alpha_1v_1+\alpha_2v_2+\cdots+\alpha_nv_n=0,
\]
to elementy są liniowo zależne.

W praktyce problem sprowadza się więc do rozwiązania jednorodnego równania liniowego lub przeanalizowania macierzy zbudowanej z badanych elementów.

\section{Metoda 1: badanie kombinacji liniowej}

Najbardziej ogólna metoda polega na zapisaniu równania
\[
\alpha_1v_1+\cdots+\alpha_nv_n=0
\]
i przekształceniu go do układu równań.

\subsection{Interpretacja wyniku}
\begin{itemize}
    \item Jeżeli jedynym rozwiązaniem jest $\alpha_1=\cdots=\alpha_n=0$, elementy są liniowo niezależne.
    \item Jeżeli układ ma rozwiązania różne od zerowego, elementy są liniowo zależne.
    \item Każde niezerowe rozwiązanie określa zależność liniową między elementami.
\end{itemize}

\subsection{Przykład -- wektory w $\mathbb R^3$}
Sprawdzić liniową niezależność wektorów
\[
v_1=\begin{bmatrix}1\\2\\3\end{bmatrix},
\qquad
v_2=\begin{bmatrix}2\\1\\0\end{bmatrix},
\qquad
v_3=\begin{bmatrix}0\\3\\6\end{bmatrix}.
\]

Równanie
\[
\alpha_1v_1+\alpha_2v_2+\alpha_3v_3=0
\]
prowadzi do układu
\[
\begin{cases}
\alpha_1+2\alpha_2=0,\\
2\alpha_1+\alpha_2+3\alpha_3=0,\\
3\alpha_1+6\alpha_3=0.
\end{cases}
\]
Z drugiego i trzeciego równania można otrzymać
\[
\alpha_1=-2\alpha_2,
\qquad
3(-2\alpha_2)+6\alpha_3=0
\quad\Longrightarrow\quad
\alpha_3=\alpha_2.
\]
Po podstawieniu do drugiego równania:
\[
2(-2\alpha_2)+\alpha_2+3\alpha_2=0.
\]
Zachodzi ono dla dowolnego $\alpha_2$, na przykład $\alpha_2=1$. Wtedy
\[
\alpha_1=-2,\qquad \alpha_3=1.
\]
Zatem
\[
-2v_1+v_2+v_3=0,
\]
więc wektory są liniowo zależne.

\section{Metoda 2: macierz i rząd}

Jeżeli $v_1,\ldots,v_n\in\mathbb K^m$, zapisuje się te wektory jako kolumny macierzy:
\[
A=\begin{bmatrix}
v_1&v_2&\cdots&v_n
\end{bmatrix}.
\]

Równanie
\[
\alpha_1v_1+\cdots+\alpha_nv_n=0
\]
jest wtedy układem jednorodnym
\[
A\alpha=0.
\]

\subsection{Kryterium rzędu}
Wektory $v_1,\ldots,v_n$ są liniowo niezależne wtedy i tylko wtedy, gdy
\[
\operatorname{rz}A=n.
\]
W szczególności:
\begin{itemize}
    \item jeżeli $n>m$, wektory w $\mathbb K^m$ są zawsze liniowo zależne;
    \item jeżeli $n=m$, macierz $A$ jest kwadratowa i wystarczy sprawdzić jej wyznacznik;
    \item różne przestrzenie liniowe wymagają zapisania ich elementów jako współrzędnych względem ustalonej bazy.
\end{itemize}

\subsection{Przykład -- rząd macierzy}
Sprawdzamy wektory
\[
v_1=\begin{bmatrix}1\\1\\0\end{bmatrix},
\qquad
v_2=\begin{bmatrix}1\\0\\1\end{bmatrix},
\qquad
v_3=\begin{bmatrix}0\\1\\-1\end{bmatrix}.
\]
Macierz zbudowana z ich współrzędnych ma postać
\[
A=
\begin{bmatrix}
1&1&0\\
1&0&1\\
0&1&-1
\end{bmatrix}.
\]
Po przeliczeniu
\[
\det A
=
1(0\cdot(-1)-1\cdot1)
-1(1\cdot(-1)-1\cdot0)
=-1+1=0.
\]
Ponieważ $\det A=0$, rząd macierzy jest mniejszy niż $3$. Wektory są liniowo zależne.

\section{Metoda 3: wyznacznik}

Jeżeli badamy dokładnie $n$ wektorów w przestrzeni $\mathbb K^n$, to można utworzyć z nich macierz kwadratową. Wektory są liniowo niezależne wtedy i tylko wtedy, gdy
\[
\det A\neq0.
\]

W przeciwnym razie:
\[
\det A=0
\quad\Longrightarrow\quad
\text{wektory są liniowo zależne}.
\]

Metoda jest szybka dla małych macierzy i ma bezpośrednią interpretację geometryczną: zerowy wyznacznik oznacza, że wektory nie rozpinają całej przestrzeni o wymiarze $n$.

\section{Metoda eliminacji Gaussa}

Można utworzyć macierz, której wierszami są współrzędne badanych elementów, i sprowadzić ją do postaci schodkowej.

Liczba niezerowych wierszy po eliminacji jest równa rzędowi macierzy. Jeżeli liczba ta jest równa liczbie badanych elementów, elementy są liniowo niezależne.

Przy tej metodzie można jednocześnie:
\begin{itemize}
    \item sprawdzić liniową niezależność,
    \item znaleźć zależności między elementami,
    \item wybrać bazę spośród podanego zbioru,
    \item sprawdzić, czy badane elementy generują całą przestrzeń.
\end{itemize}

\section{Metoda macierzy Grama}

W przypadku przestrzeni z iloczynem skalarnym można utworzyć macierz Grama:
\[
G=
\begin{bmatrix}
\langle v_1,v_1\rangle & \langle v_1,v_2\rangle & \cdots & \langle v_1,v_n\rangle\\
\langle v_2,v_1\rangle & \langle v_2,v_2\rangle & \cdots & \langle v_2,v_n\rangle\\
\vdots & \vdots & \ddots & \vdots\\
\langle v_n,v_1\rangle & \langle v_n,v_2\rangle & \cdots & \langle v_n,v_n\rangle
\end{bmatrix}.
\]

Wektory $v_1,\ldots,v_n$ są liniowo niezależne wtedy i tylko wtedy, gdy
\[
\det G>0.
\]
Dla układu zespolonego macierz Grama jest dodatnio określona dokładnie w przypadku liniowej niezależności.

Praktycznie metodę tę stosuje się szczególnie wtedy, gdy elementami są funkcje, ciągi lub inne obiekty, dla których iloczyn skalarny jest znany, a zapis macierzowy współrzędnych może być trudny.

\subsection{Wniosek dla układu ortogonalnego}
Jeżeli wszystkie badane niezerowe elementy są parami ortogonalne:
\[
\langle v_i,v_j\rangle=0 \quad (i\neq j),
\]
to tworzą układ liniowo niezależny.

Istotnie, z
\[
\alpha_1v_1+\cdots+\alpha_nv_n=0
\]
wynika po wykonaniu iloczynu skalarnego z $v_i$:
\[
\alpha_i\langle v_i,v_i\rangle=0.
\]
Ponieważ $v_i\neq0$, mamy $\langle v_i,v_i\rangle>0$, a więc $\alpha_i=0$ dla każdego $i$.

\section{Metoda sprawdzania wielomianów i funkcji}

Dla wielomianów najczęściej porównuje się współczynniki przy jednakowych potęgach. Należy zauważyć, że wielomiany stopnia co najwyżej $n$ mają naturalną bazę
\[
1,x,x^2,\ldots,x^n.
\]
Dlatego można przedstawić wielomiany jako wektory współczynników i zastosować rząd lub wyznacznik.

Na przykład wielomiany
\[
p_1(x)=1+x,\qquad p_2(x)=1-x,\qquad p_3(x)=x+x^2
\]
mają współrzędne:
\[
p_1=(1,1,0),\qquad p_2=(1,-1,0),\qquad p_3=(0,1,1)
\]
w bazie $(1,x,x^2)$.
Z postaci macierzy
\[
A=
\begin{bmatrix}
1&1&0\\
1&-1&1\\
0&0&1
\end{bmatrix}
\]
wynika $\det A=-2\neq0$, więc wielomiany są liniowo niezależne.

\subsection{Dla funkcji: badanie tożsamości}
Jeżeli badamy funkcje $f_1,\ldots,f_n$, zapisujemy
\[
\alpha_1f_1(t)+\cdots+\alpha_nf_nf(t)=0
\]
dla wszystkich $t$ z dziedziny i podstawiamy różne wartości $t$ lub różnicujemy, jeżeli funkcje są odpowiednio gładkie.

Na przykład układ $(1,t,e^t)$ można badać, podstawiając kilka wartości $t$. Wyznaczenie współczynników z układu równań pozwala sprawdzić, czy istnieje niezerowa kombinacja zerowa.

\section{Wyznacznik Wrońskiego}

Dla $n$ funkcji różniczkowalnych $f_1,\ldots,f_n$ można czasem użyć wyznacznika Wrońskiego:
\[
W(t)=
\det
\begin{bmatrix}
f_1(t)&f_2(t)&\cdots&f_n(t)\\
f_1'(t)&f_2'(t)&\cdots&f_n'(t)\\
\vdots&\vdots&\ddots&\vdots\\
f_1^{(n-1)}(t)&f_2^{(n-1)}(t)&\cdots&f_n^{(n-1)}(t)
\end{bmatrix}.
\]

\begin{itemize}
    \item Jeśli $W(t_0)\neq0$ dla pewnego punktu $t_0$, funkcje są liniowo niezależne.
    \item Jeśli $W(t)=0$ dla wszystkich $t$, nie można na tej podstawie samodzielnie wnioskować o zależności liniowej.
\end{itemize}

Ważne: zerowy Wrońskian nie zawsze oznacza liniową zależność. Wykorzystanie tej metody wymaga ostrożności i znajomości dodatkowych założeń, na przykład o funkcjach będących rozwiązaniami pewnego równania różniczkowego liniowego.

\section{Najczęstsze metody zależnie od rodzaju zadania}

\begin{center}
\begin{tabular}{|p{0.32\linewidth}|p{0.56\linewidth}|}
\hline
\textbf{Rodzaj elementów} & \textbf{Najdogodniejsza metoda} \\
\hline
Wektory z $\mathbb R^n$ lub $\mathbb C^n$ & Układ jednorodny, rząd macierzy, wyznacznik \\
\hline
Dokładnie $n$ wektorów w $n$ wymiarach & Wyznacznik macierzy kwadratowej \\
\hline
Więcej wektorów niż wymiar & Od razu liniowa zależność \\
\hline
Wielomiany & Współczynniki względem bazy jednomianów \\
\hline
Funkcje & Układ wartości, pochodne, macierz Grama lub Wrońskian \\
\hline
Układ ortogonalny & Iloczyn skalarny z każdym elementem układu \\
\hline
Duże układy wektorów & Eliminacja Gaussa i rząd macierzy \\
\hline
\end{tabular}
\end{center}

\section{Podsumowanie algorytmu}
\begin{enumerate}
    \item Zapisz równanie kombinacji liniowej równej elementowi zerowemu.
    \item Przedstaw elementy za pomocą współrzędnych, macierzy albo układu równań.
    \item Znajdź rozwiązania jednorodne lub policz rząd macierzy.
    \item Jeżeli rozwiązaniem jest tylko wektor zerowy albo rząd równy jest liczbie elementów, układ jest liniowo niezależny.
    \item Jeżeli istnieje niezerowe rozwiązanie albo rząd jest mniejszy niż liczba elementów, układ jest liniowo zależny.
\end{enumerate}

\end{document}
